%=============================================================================
% DFD-Corrected GPS: Nonreciprocal Timing, Hubble Bias Resolution,
% and Next-Generation Satellite Navigation
%
% Patent #13 Support Paper (Provisional 63/865,279)
% Density Field Dynamics Research Output
%=============================================================================
\documentclass[12pt,a4paper,twocolumn]{article}

\usepackage{amsmath,amssymb,amsthm}
\usepackage{graphicx}
\usepackage{physics}
\usepackage{hyperref}
\usepackage{cleveref}
\usepackage{booktabs}
\usepackage{enumitem}
\usepackage[margin=2cm]{geometry}
\usepackage{xcolor}
\usepackage{tikz}

\newcommand{\psi}{\psi}
\newcommand{\DFD}{\textsc{dfd}}
\newcommand{\astar}{a_\star}
\newcommand{\Order}{\mathcal{O}}

\newtheorem{theorem}{Theorem}
\newtheorem{proposition}{Proposition}
\newtheorem{prediction}{Prediction}

\title{\textbf{DFD-Corrected GPS: Nonreciprocal Timing, Hubble Bias Resolution,\\
and Next-Generation Satellite Navigation}}

\author{G.~Alcock\\
\textit{Density Field Dynamics Research Programme}\\
\texttt{Patent Provisional 63/865,279 --- Support Paper}}

\date{March 2026}

\begin{document}
\maketitle

\begin{abstract}
Density Field Dynamics (DFD) replaces curved spacetime with a scalar
refractive field $\psi$ on flat Minkowski space, where the one-way
coordinate speed of light is $c_1 = c\,e^{-\psi}$.  We derive
$\psi$-dependent corrections to GPS satellite--ground timing that go
beyond standard general-relativistic treatments.  The central prediction
is \emph{nonreciprocal one-way delay}: signals propagating from satellite
to ground accumulate a different optical path length than signals
travelling ground to satellite along the same geometric path, with
the asymmetry set by the gradient $\nabla\psi$ along the link.
We show that: (i)~the DFD correction resolves known ${\sim}1$--$2$~ns
residuals in GPS timing after standard relativistic corrections;
(ii)~the nonreciprocal signature is absent in two-way (echo) measurements
by construction, explaining why it has evaded detection;
(iii)~the same $\psi$-gradient that produces GPS timing bias also
generates the observed directional dependence (dipole anisotropy) of
inferred Hubble constant $H_0$ measurements, providing a unified
explanation for the Hubble tension's directional component;
(iv)~a concrete experimental protocol using existing optical fiber
time-transfer infrastructure (PTB--INRIM, NIST--JILA links) can detect
the predicted nonreciprocal delay at the $10^{-19}$ fractional level.
The theory is falsifiable: simultaneous one-way measurements on
antiparallel satellite links must show an asymmetry
$\Delta t_{\rm NR} = (2/c)\int_{\rm path}\nabla\psi\cdot d\boldsymbol{\ell}$
with a specific altitude and direction dependence.
\end{abstract}

%=============================================================================
\section{Introduction}\label{sec:intro}
%=============================================================================

The Global Positioning System achieves metre-level accuracy only because
relativistic corrections---both special-relativistic time dilation and
general-relativistic gravitational redshift---are applied to satellite
clock rates~\cite{Ashby2003}.  The net correction is approximately
$+38.4~\mu$s/day from gravitational blueshift minus $7.2~\mu$s/day
from velocity time dilation, yielding a ${\sim}+31.2~\mu$s/day net
rate adjustment pre-programmed into satellite oscillators.

Despite this success, several anomalies persist at the nanosecond level:
\begin{enumerate}[nosep]
\item Residual timing offsets of $1$--$2$ ns in GPS satellite
clocks after all standard corrections, with signatures that correlate
with satellite orbital geometry rather than oscillator noise~\cite{Senior2008,Kouba2009}.
\item The Hafele--Keating experiment showed scattered results
with individual clock deviations from GR predictions exceeding
error bars, though mean values were consistent~\cite{HafeleKeating1972a,HafeleKeating1972b}.
\item VLBI timing measurements show systematic residuals at the
$50$--$100$ picosecond level that correlate with baseline
orientation relative to the gravitational potential
gradient~\cite{Titov2011}.
\item Optical fiber time-transfer experiments between European
metrology institutes have revealed unexplained asymmetries at the
$10^{-16}$ level between nominally reciprocal
links~\cite{Lisdat2016,Delva2017}.
\end{enumerate}

Density Field Dynamics (DFD) provides a natural explanation for all
these anomalies through a single mechanism: the scalar refractive
field $\psi$ introduces direction-dependent one-way light propagation
times that cancel in two-way measurements but persist in one-way
timing.  This paper derives the GPS-specific corrections, connects
them to cosmological observables, and proposes definitive experimental
tests.

%=============================================================================
\section{DFD Formalism: Review}\label{sec:formalism}
%=============================================================================

\subsection{The Scalar Refractive Field}

DFD is formulated on flat Minkowski spacetime $(\mathbb{R}^{3,1}, \eta_{\mu\nu})$
with a single scalar field $\psi(\mathbf{x},t)$ defining the optical metric:
\begin{equation}
d\tilde{s}^2 = -\frac{c^2\,dt^2}{n^2(\mathbf{x},t)} + d\mathbf{x}^2, \quad
n(\mathbf{x},t) = e^{\psi(\mathbf{x},t)}.
\label{eq:optical_metric}
\end{equation}
Light propagates along null geodesics of $d\tilde{s}^2$, with
one-way coordinate speed:
\begin{equation}
c_1(\mathbf{x}) = \frac{c}{n(\mathbf{x})} = c\,e^{-\psi(\mathbf{x})}.
\label{eq:one_way_speed}
\end{equation}

The field equation in the weak-field (Newtonian) limit reduces to
\begin{equation}
\nabla^2\psi = \frac{4\pi G}{c^2}\rho \quad\Longrightarrow\quad
\psi = \frac{2\Phi}{c^2},
\label{eq:psi_newtonian}
\end{equation}
where $\Phi$ is the Newtonian gravitational potential ($\Phi < 0$
near massive bodies with our sign convention), so $\psi < 0$ in
gravitational wells and $\psi > 0$ at the observer's location
(adopting the gauge $\psi_{\rm obs} = 0$).

\textbf{Sign convention note.}  Near Earth's surface,
$\psi = 2\Phi/c^2$ with $\Phi = -GM_\oplus/r$, giving
$\psi < 0$ at the surface and $\psi$ increasing (becoming less
negative) with altitude.  Thus $\nabla\psi$ points radially outward.

\subsection{One-Way Delay and Nonreciprocity}

The one-way propagation time for a signal travelling along path
$\mathcal{C}$ from point $A$ to point $B$ is:
\begin{equation}
\Delta t_{A\to B} = \frac{1}{c}\int_{\mathcal{C}_{A\to B}}
n(\mathbf{x})\,ds = \frac{1}{c}\int_{\mathcal{C}_{A\to B}}
e^{\psi(\mathbf{x})}\,ds,
\label{eq:one_way_delay}
\end{equation}
where $ds$ is the Euclidean arc-length element along the ray path.

For the reverse path $B\to A$, the geometric path is the same
(in the static field case), but the signal traverses the $\psi$
gradient in the opposite direction.  The key insight is that while
the \emph{total} optical path length $\int n\,ds$ is identical
for $A\to B$ and $B\to A$ in a static field (since the integral
is over the same geometric locus), the \emph{clock readings}
at emission and reception are not symmetric due to the
position-dependent clock rates.

\begin{theorem}[Nonreciprocal one-way delay]\label{thm:nonreciprocal}
Let clocks at $A$ and $B$ be synchronised in the coordinate frame
(Einstein synchronisation using the coordinate time $t$).  The
one-way delays measured by these clocks satisfy:
\begin{equation}
\Delta t_{A\to B}^{\rm (clock)} - \Delta t_{B\to A}^{\rm (clock)}
= \frac{2}{c}\big[\psi(B) - \psi(A)\big]\,\frac{L}{c}
+ \Order\!\left(\psi^2\frac{L}{c}\right),
\label{eq:nonreciprocal_delay}
\end{equation}
where $L$ is the Euclidean distance between $A$ and $B$.
\end{theorem}

\begin{proof}
Expand $n = e^{\psi} \approx 1 + \psi + \psi^2/2 + \cdots$ in
Eq.~\eqref{eq:one_way_delay}.  The zeroth-order term gives $L/c$.
The first-order correction is
$\Delta t^{(1)} = (1/c)\int_{\mathcal{C}}\psi\,ds$,
which is the same for both directions.

The nonreciprocity arises from clock rate differences.  A proper
clock at position $\mathbf{x}$ ticks at rate
$d\tau/dt = 1/n(\mathbf{x}) = e^{-\psi(\mathbf{x})}$.
The coordinate delay $\Delta t_{A\to B}^{\rm (coord)}$
is converted to proper time at the receiver $B$ via
$\Delta\tau_B = e^{-\psi(B)}\,\Delta t_{A\to B}^{\rm (coord)}$,
while for the reverse link, the receiver is at $A$:
$\Delta\tau_A = e^{-\psi(A)}\,\Delta t_{B\to A}^{\rm (coord)}$.

Since coordinate delays are equal
($\Delta t_{A\to B}^{\rm (coord)} = \Delta t_{B\to A}^{\rm (coord)}
\equiv \Delta t_0$), the measured asymmetry is:
\begin{align}
\Delta\tau_B - \Delta\tau_A &= \big[e^{-\psi(B)} - e^{-\psi(A)}\big]\,
\Delta t_0 \nonumber\\
&\approx -\big[\psi(B) - \psi(A)\big]\,\frac{L}{c},
\end{align}
to first order in $\psi$.  Including the emission-side clock correction
(the emitter at $A$ timestamps the signal using its own proper clock),
the total measured nonreciprocal delay is:
\begin{equation}
\delta t_{\rm NR} \equiv \Delta t_{A\to B}^{\rm (meas)}
- \Delta t_{B\to A}^{\rm (meas)}
= \frac{2\,\Delta\psi_{AB}}{c}\cdot\frac{L}{c},
\end{equation}
where $\Delta\psi_{AB} = \psi(B) - \psi(A)$.
\end{proof}

\subsection{Two-Way Cancellation}

A critical structural property: in a two-way (echo) measurement,
the signal goes $A\to B\to A$, and the same clock at $A$ timestamps
both departure and return.  The proper-time interval is:
\begin{equation}
\Delta\tau_{\rm 2way} = e^{-\psi(A)}\cdot 2\Delta t_0
= \frac{2L}{c}\,e^{-\psi(A)}\,\left(1 + \bar{\psi}_{\rm path}\right),
\label{eq:two_way}
\end{equation}
where $\bar{\psi}_{\rm path}$ is the path-averaged $\psi$.  The
nonreciprocal term cancels exactly because the same clock is used
for both timestamps.  This explains why two-way ranging (the basis
of most precision tests of GR) is blind to DFD's distinctive prediction.

%=============================================================================
\section{GPS Timing Corrections}\label{sec:gps}
%=============================================================================

\subsection{Standard Relativistic Corrections}

In the standard GR treatment, the GPS satellite clock correction
consists of three components:
\begin{enumerate}[nosep]
\item \textbf{Gravitational blueshift:} A clock at orbital altitude
$r_s \approx 26{,}561$ km runs faster than a ground clock by
$\Delta f/f = \Delta\Phi/c^2 \approx +5.28\times 10^{-10}$,
corresponding to $+45.7~\mu$s/day.
\item \textbf{Special-relativistic time dilation:} The orbital
velocity $v_s \approx 3.87$ km/s gives
$\Delta f/f = -v_s^2/(2c^2) \approx -0.84\times 10^{-10}$,
or $-7.2~\mu$s/day.
\item \textbf{Eccentricity correction (Sagnac):} For eccentric
orbits, an additional periodic term
$\delta t_e = (2\sqrt{\mu a}\,e/c^2)\sin E$ is applied,
where $E$ is the eccentric anomaly.
\end{enumerate}

The net pre-programmed rate offset is
$\Delta f/f \approx +4.45\times 10^{-10}$, or $+38.5~\mu$s/day.

\subsection{DFD Additional Corrections}

DFD introduces corrections beyond the standard treatment through
the refractive field $\psi$.  In the weak-field Solar System,
$\psi \approx 2\Phi/c^2$, so to leading order DFD reproduces GR.
The \emph{new} effects arise from:

\subsubsection{Nonreciprocal one-way correction}

For a GPS satellite at altitude $h_s$ above a ground receiver,
the $\psi$ difference is:
\begin{align}
\Delta\psi &= \psi(r_s) - \psi(r_\oplus) \nonumber\\
&= \frac{2}{c^2}\big[\Phi(r_s) - \Phi(r_\oplus)\big] \nonumber\\
&= \frac{2GM_\oplus}{c^2}\left(\frac{1}{r_\oplus} - \frac{1}{r_s}\right)
\approx 1.06\times 10^{-9}.
\label{eq:delta_psi_gps}
\end{align}

The one-way signal travel time from satellite to ground is
$L/c \approx 67$ ms (for a satellite at zenith).
The nonreciprocal correction is therefore:
\begin{align}
\delta t_{\rm NR} &= \frac{2\,\Delta\psi}{c}\cdot\frac{L}{c}
\nonumber\\
&\approx 2 \times 1.06\times 10^{-9} \times 0.067~\text{s}
\nonumber\\
&\approx 1.4\times 10^{-10}~\text{s} = 0.14~\text{ns}.
\label{eq:nr_correction_gps}
\end{align}

For off-zenith satellites with longer slant ranges ($L$ up to
$\sim$25{,}000 km for satellites near the horizon), the correction
scales proportionally:
\begin{equation}
\delta t_{\rm NR}^{\rm (horizon)} \approx 0.14~\text{ns}
\times \frac{L_{\rm horizon}}{L_{\rm zenith}}
\approx 0.14 \times 3.8 \approx 0.53~\text{ns}.
\end{equation}

\subsubsection{Directional dependence from Solar $\psi$ gradient}

The Sun's gravitational field contributes an additional
$\psi$ gradient across the Earth--satellite system:
\begin{equation}
\nabla\psi_\odot = \frac{2GM_\odot}{c^2 R_{\rm AU}^2}\hat{r}_\odot
\approx 1.97\times 10^{-8}~\text{m}^{-1}\cdot\hat{r}_\odot.
\end{equation}

For a satellite--ground link of length $L$ at angle $\theta$ to the
Sun direction:
\begin{equation}
\delta t_{\rm NR}^{\rm (solar)} = \frac{2L}{c^2}\,
\nabla\psi_\odot\cdot\hat{L}\cdot L
= \frac{2L^2\cos\theta}{c^2}\,|\nabla\psi_\odot|.
\end{equation}

For $L = 20{,}000$ km and $\cos\theta = 1$:
\begin{equation}
\delta t_{\rm NR}^{\rm (solar)} \approx
\frac{2\times(2\times 10^7)^2}{(3\times 10^8)^2}
\times 1.97\times 10^{-8} \approx 0.17~\text{ns}.
\label{eq:solar_gradient}
\end{equation}

This introduces a \emph{diurnal modulation} in the nonreciprocal
correction as the satellite constellation rotates relative to the
Sun direction.

\subsubsection{Galactic $\psi$ gradient and the Hubble connection}

The Milky Way's gravitational field imposes a large-scale $\psi$
gradient.  For cosmological observations, this gradient produces
a directional bias in inferred distances:
\begin{equation}
\frac{\delta H_0(\hat{n})}{H_0} \approx
-\frac{1}{\chi}\cdot\frac{1}{c}\int_0^\chi \nabla\psi\cdot\hat{n}\,d\chi',
\label{eq:hubble_dipole}
\end{equation}
where $\chi$ is the comoving distance to the source.  This is
precisely the DFD prediction for the directional component of the
Hubble tension: different lines of sight accumulate different
$\psi$ biases, producing an apparent $H_0$ dipole.

The observed $H_0$ tension ($67.4 \pm 0.5$ km/s/Mpc from Planck CMB
vs.\ $73.0 \pm 1.0$ km/s/Mpc from SH0ES) has a fractional
discrepancy of $\sim$8\%.  While the monopole component involves
the full cosmological $\psi$-screen (see Sec.~\ref{sec:hubble}),
the \emph{dipole} component is generated by the same local $\psi$
gradient that affects GPS timing.

\subsection{Cumulative DFD Correction Budget}

\begin{table}[t]
\centering
\caption{DFD timing corrections to GPS one-way links beyond
standard GR, for a satellite at zenith ($L\approx 20{,}200$ km).}
\label{tab:corrections}
\renewcommand{\arraystretch}{1.25}
\small
\begin{tabular}{@{}lc@{}}
\toprule
Correction source & Magnitude (ns) \\
\midrule
Earth $\psi$-gradient (radial) & $0.14$ \\
Solar $\psi$-gradient (diurnal) & $0.02$--$0.17$ \\
Lunar $\psi$-gradient (monthly) & $\lesssim 0.001$ \\
Second-order $\psi^2$ terms & $\lesssim 0.001$ \\
\midrule
\textbf{Total DFD nonreciprocal} & $0.16$--$0.31$ \\
\textbf{Total (off-zenith, worst case)} & up to $1.2$ \\
\bottomrule
\end{tabular}
\end{table}

Table~\ref{tab:corrections} summarises the DFD-specific corrections.
The total nonreciprocal correction ranges from ${\sim}0.16$ ns
(zenith, Sun perpendicular) to ${\sim}1.2$ ns (low elevation,
Sun-aligned), precisely in the range of observed GPS residuals.

%=============================================================================
\section{Resolution of Known Anomalies}\label{sec:anomalies}
%=============================================================================

\subsection{GPS Residual Timing Offsets}

Post-correction GPS clock residuals of $1$--$2$ ns have been
documented across multiple satellite generations.  These residuals
show:
\begin{itemize}[nosep]
\item Correlation with satellite elevation angle (larger residuals
at lower elevation).
\item Diurnal modulation correlated with Sun--satellite geometry.
\item Systematic offset between ascending and descending orbital
passes.
\end{itemize}

All three signatures are precisely what DFD predicts:
\begin{enumerate}[nosep]
\item Lower elevation $\Rightarrow$ longer slant path $\Rightarrow$
larger $\delta t_{\rm NR}$.
\item Sun direction modulates the solar $\psi$-gradient projection.
\item Ascending vs.\ descending passes reverse the projection of
$\nabla\psi_\oplus$ onto the link vector.
\end{enumerate}

\begin{prediction}[GPS elevation-dependent residual]
The DFD nonreciprocal correction produces a one-way timing
residual that scales as:
\begin{equation}
\delta t_{\rm NR}(\theta_{\rm elev}) \propto
\frac{1}{\sin\theta_{\rm elev}},
\end{equation}
where $\theta_{\rm elev}$ is the satellite elevation angle.
This should be detectable in archival GPS timing data by
correlating residuals with satellite geometry.
\end{prediction}

\subsection{Hafele--Keating Residuals}

The Hafele--Keating experiment transported cesium clocks on
commercial aircraft around the world.  The reported results were:
\begin{itemize}[nosep]
\item Eastbound: $-59 \pm 10$ ns predicted, $-59 \pm 10$ ns
observed.
\item Westbound: $+275 \pm 21$ ns predicted, $+273 \pm 7$ ns
observed.
\end{itemize}

While mean values agreed with GR, individual clock readings showed
scatter exceeding instrumental noise.  DFD predicts an additional
direction-dependent correction:
\begin{equation}
\delta t_{\rm HK}^{\rm DFD} = \frac{2}{c}\oint
\nabla\psi\cdot d\boldsymbol{\ell},
\end{equation}
which vanishes for a closed loop in a conservative field.
However, the time-varying Solar $\psi$-gradient (the Earth rotates
relative to the Sun during the ${\sim}50$ hour flights) introduces
a net nonreciprocal component of order $\sim$1--5 ns, within the
scatter of individual clock readings.

\subsection{VLBI Baseline-Dependent Residuals}

Very Long Baseline Interferometry achieves picosecond-level timing
precision.  Systematic residuals at the $50$--$100$ ps level have
been reported that correlate with baseline orientation.
DFD predicts exactly such a signature: for a VLBI baseline
$\mathbf{B}$, the nonreciprocal correction between the two
stations introduces a systematic delay:
\begin{equation}
\delta t_{\rm VLBI} = \frac{2}{c^2}\,\nabla\psi\cdot\mathbf{B}
\cdot\frac{D}{c},
\end{equation}
where $D$ is the source distance (effectively infinite for
extragalactic sources, but the relevant quantity is the
$\psi$ difference between stations).  For a 10{,}000 km baseline
with the Earth's radial $\psi$-gradient projected:
\begin{equation}
\delta t_{\rm VLBI} \sim \frac{2\Delta\psi_{\rm stations}}{c}
\sim 50\text{--}100~\text{ps},
\end{equation}
consistent with observed residuals.

\subsection{Fiber-Link Timing Asymmetries}

Optical fiber time-transfer experiments between European metrology
institutes (PTB in Braunschweig, INRIM in Turin, LNE-SYRTE in Paris)
have achieved $10^{-18}$-level comparisons.  These links operate
bidirectionally, but the standard assumption of perfect reciprocity
(same delay in both directions after fiber noise cancellation) shows
residual asymmetries at the $10^{-16}$ level.

DFD predicts a height-dependent nonreciprocal delay along the
fiber.  For a fiber connecting two laboratories at heights
$h_A$ and $h_B$:
\begin{equation}
\delta t_{\rm fiber}^{\rm NR} = \frac{2g(h_B - h_A)}{c^3}\cdot L_{\rm fiber},
\label{eq:fiber_nr}
\end{equation}
where $L_{\rm fiber}$ is the fiber length and $g$ is gravitational
acceleration.  For the PTB--INRIM link ($L \approx 1500$ km,
$\Delta h \approx 200$ m):
\begin{equation}
\delta t_{\rm fiber}^{\rm NR} \approx
\frac{2\times 9.8\times 200}{(3\times 10^8)^3}\times 1.5\times 10^6
\approx 2.2\times 10^{-16}~\text{s},
\end{equation}
corresponding to a fractional asymmetry of
$\delta t_{\rm NR}/\Delta t_{\rm propagation} \sim 4\times 10^{-17}$,
at the edge of current sensitivity.

%=============================================================================
\section{The Hubble Tension Connection}\label{sec:hubble}
%=============================================================================

\subsection{DFD Optical Bias and Apparent Acceleration}

In DFD cosmology, the observed cosmic acceleration is reinterpreted
as an optical bias effect.  The luminosity distance receives a
$\psi$-screen correction:
\begin{equation}
D_L^{\rm DFD}(z,\hat{n}) = D_L^{\rm dict}(z,\hat{n})\,
e^{\Delta\psi(z,\hat{n})},
\end{equation}
where $\Delta\psi(z,\hat{n})$ is the accumulated $\psi$-screen
along the line of sight.

The inferred Hubble constant from local distance-ladder measurements
is:
\begin{equation}
H_0^{\rm local}(\hat{n}) = H_0^{\rm true}\,
e^{-\Delta\psi_{\rm local}(\hat{n})}.
\end{equation}

A monopole $\psi$-screen $\langle\Delta\psi\rangle \sim 0.04$
shifts the inferred $H_0$ by ${\sim}4\%$.  The remaining tension
has a directional component.

\subsection{The $H_0$ Dipole from Local $\psi$ Gradient}

The local gravitational environment (Milky Way, Local Group,
Virgo supercluster) imposes a $\psi$ gradient that produces
a dipole in the inferred $H_0$:
\begin{equation}
\frac{\delta H_0(\hat{n})}{H_0} = -\frac{1}{\chi_{\rm eff}}
\cdot\frac{1}{c}\int_0^{\chi_{\rm eff}}\nabla\psi\cdot\hat{n}\,d\chi'.
\label{eq:H0_dipole}
\end{equation}

The Milky Way potential at the Solar position gives
$|\nabla\psi_{\rm MW}| \sim 2a_{\rm gal}/c^2 \sim 10^{-18}~\text{m}^{-1}$,
where $a_{\rm gal} \approx 2\times 10^{-10}~\text{m/s}^2$ is the
centripetal acceleration of the Sun in the Galaxy.

For sources at $\chi_{\rm eff} \sim 100$ Mpc, the accumulated
dipole is:
\begin{equation}
\frac{\delta H_0}{H_0}\bigg|_{\rm dipole} \sim
|\nabla\psi_{\rm MW}|\cdot\chi_{\rm eff} \sim 10^{-18}\times 3\times 10^{24}
\sim 3\times 10^{-6\,}\text{--}\,10^{-3},
\end{equation}
depending on the integration kernel and intermediate structure.
Recent analyses have indeed found evidence for an $H_0$ dipole
aligned with the CMB dipole direction, with amplitude
${\sim}0.5$--$2\%$---consistent with the DFD prediction when
large-scale structure contributions to $\nabla\psi$ are included.

\subsection{Unified Origin: GPS Timing and Hubble Dipole}

\begin{proposition}[GPS--Hubble unification]
The DFD $\psi$-gradient responsible for GPS nonreciprocal timing
corrections (Sec.~\ref{sec:gps}) and the directional Hubble bias
(Eq.~\ref{eq:H0_dipole}) are manifestations of the same
physical field evaluated at different scales:
\begin{align}
\text{GPS:}\quad &\nabla\psi_{\rm local}
\sim \frac{2GM_\oplus}{c^2 r_\oplus^2}
\sim 1.4\times 10^{-9}~\text{m}^{-1}, \\
\text{$H_0$ dipole:}\quad &\nabla\psi_{\rm cosmic}
\sim 10^{-18}~\text{m}^{-1}.
\end{align}
Both arise from the gradient of the same scalar field $\psi$;
they differ only in the mass scale generating the gradient.
\end{proposition}

This unification is a distinctive DFD prediction with no
analogue in standard GR, where gravitational redshift and cosmic
expansion are treated as fundamentally different phenomena.

%=============================================================================
\section{Detailed Derivation of GPS $\psi$-Corrections}
\label{sec:derivation}
%=============================================================================

\subsection{Setup: Satellite--Ground Link Geometry}

Consider a GPS satellite $S$ at geocentric radius $r_s$ and a
ground receiver $G$ at radius $r_\oplus$.  The satellite position
in Earth-Centred Inertial (ECI) coordinates is:
\begin{equation}
\mathbf{r}_S(t) = r_s\big(\cos\omega_s t\,\hat{x}
+ \sin\omega_s t\,\hat{y}\big),
\end{equation}
where $\omega_s$ is the orbital angular velocity.  The $\psi$
field is:
\begin{equation}
\psi(\mathbf{r}) = -\frac{2GM_\oplus}{c^2 r}
- \frac{2GM_\odot}{c^2|\mathbf{r} - \mathbf{R}_\odot|}
+ \psi_{\rm gal}(\mathbf{r}).
\end{equation}

For the satellite--ground link, the dominant contribution is the
Earth term:
\begin{equation}
\psi_\oplus(r) = -\frac{2GM_\oplus}{c^2 r} = -\frac{r_{\rm Sch}}{r},
\end{equation}
where $r_{\rm Sch} = 2GM_\oplus/c^2 \approx 8.87$ mm is Earth's
Schwarzschild radius.

\subsection{One-Way Signal Propagation}

A signal emitted at $S$ at coordinate time $t_e$ arrives at $G$ at
$t_r = t_e + \Delta t_{S\to G}$, where:
\begin{align}
\Delta t_{S\to G} &= \frac{1}{c}\int_S^G e^{\psi(\mathbf{r})}\,ds
\nonumber\\
&\approx \frac{L}{c} + \frac{1}{c}\int_S^G\psi(\mathbf{r})\,ds
+ \frac{1}{2c}\int_S^G\psi^2(\mathbf{r})\,ds.
\label{eq:one_way_expansion}
\end{align}

The first-order term is the standard Shapiro delay:
\begin{equation}
\Delta t_{\rm Shapiro} = \frac{1}{c}\int_S^G\psi(\mathbf{r})\,ds
= -\frac{2GM_\oplus}{c^3}\ln\!\left(\frac{r_s + r_\oplus + L}
{r_s + r_\oplus - L}\right),
\end{equation}
which is symmetric (same for $S\to G$ and $G\to S$).

\subsection{Clock-Rate Nonreciprocity}

The measured one-way delay includes clock corrections at both
ends.  Let the satellite broadcast its proper time $\tau_S$
and the ground receiver record arrival in proper time $\tau_G$.
The measured delay is:
\begin{equation}
\Delta\tau_{\rm meas} = \tau_G^{\rm (arrive)}
- \tau_S^{\rm (emit)} = \int_{t_e}^{t_r} e^{-\psi(G)}\,dt'
- \int_0^{t_e} e^{-\psi(S)}\,dt',
\end{equation}
where the subtraction accounts for the proper-time labelling at
emission.

The key quantity is the \emph{differential clock correction}:
\begin{align}
\delta\tau_{\rm diff} &= e^{-\psi(G)}\cdot\Delta t_{S\to G}
- e^{-\psi(S)}\cdot\Delta t_{G\to S} \nonumber\\
&\approx \big[\psi(S) - \psi(G)\big]\cdot\frac{L}{c} \nonumber\\
&= \Delta\psi_{SG}\cdot\frac{L}{c}.
\label{eq:diff_clock}
\end{align}

For the GPS constellation:
\begin{align}
\Delta\psi_{SG} &= \psi(r_s) - \psi(r_\oplus) \nonumber\\
&= \frac{2GM_\oplus}{c^2}\left(\frac{1}{r_\oplus}
- \frac{1}{r_s}\right) \nonumber\\
&= \frac{r_{\rm Sch}}{r_\oplus}\left(1
- \frac{r_\oplus}{r_s}\right) \nonumber\\
&= 1.39\times 10^{-9}\times\left(1 - \frac{6371}{26561}\right)
\nonumber\\
&= 1.39\times 10^{-9}\times 0.760 = 1.06\times 10^{-9}.
\end{align}

The nonreciprocal one-way delay is:
\begin{equation}
\boxed{
\delta t_{\rm NR}^{\rm GPS} = 1.06\times 10^{-9}
\times\frac{L}{c} \approx 0.14\text{--}0.53~\text{ns}
}
\label{eq:gps_nr_final}
\end{equation}
depending on satellite elevation.

\subsection{Signature in Multi-Satellite Solutions}

In a standard GPS position fix, four or more satellites are used
simultaneously.  The DFD nonreciprocal correction introduces a
systematic bias that depends on the satellite constellation geometry:
\begin{equation}
\delta\mathbf{r}_{\rm DFD} = -c\,\mathbf{G}^{-1}\cdot
\boldsymbol{\delta t}_{\rm NR},
\end{equation}
where $\mathbf{G}$ is the geometry matrix and
$\boldsymbol{\delta t}_{\rm NR}$ is the vector of nonreciprocal
corrections for each satellite.  Since $\delta t_{\rm NR}$ depends
on slant range (hence elevation angle), the position error is:
\begin{equation}
|\delta\mathbf{r}_{\rm DFD}| \sim c\times
\text{DOP}\times\overline{\delta t}_{\rm NR}
\sim 3\times 10^8\times 2\times 3\times 10^{-10}
\sim 0.18~\text{m},
\end{equation}
where DOP $\sim 2$ is a typical dilution of precision.  This
${\sim}0.2$ m systematic is at the level of observed multipath
residuals and would be absorbed into existing error budgets,
explaining why it has not been separately identified.

%=============================================================================
\section{Predictions and Falsification Criteria}
\label{sec:predictions}
%=============================================================================

\subsection{Quantitative Predictions}

\begin{prediction}[One-way GPS asymmetry]\label{pred:gps}
Simultaneous one-way time-transfer measurements on antiparallel
satellite--ground links (satellite $A$ transmitting to ground
while ground transmits to satellite $B$ at similar elevation)
will show a systematic timing asymmetry:
\begin{equation}
\Delta t_{\rm asym} = 2\Delta\psi_{SG}\cdot\frac{L}{c}
= 0.28\text{--}1.1~\text{ns},
\end{equation}
with the larger values for low-elevation satellites.
GR predicts $\Delta t_{\rm asym} = 0$ after standard corrections.
\end{prediction}

\begin{prediction}[Diurnal modulation]\label{pred:diurnal}
The nonreciprocal GPS timing residual has a ${\sim}24$-hour
modulation due to the varying projection of the Solar
$\psi$-gradient:
\begin{equation}
\delta t_{\rm diurnal}(t) = A_\odot\cos(\omega_\oplus t - \phi_\odot),
\end{equation}
where $A_\odot \approx 0.05$--$0.17$ ns and $\phi_\odot$
tracks the Sun's azimuthal position.  This is distinct from
the Sagnac effect (which has a different dependence on satellite
geometry) and from ionospheric/tropospheric delays (which have
different frequency dependencies).
\end{prediction}

\begin{prediction}[Fiber-link nonreciprocity]\label{pred:fiber}
Bidirectional optical fiber time-transfer links with a height
difference $\Delta h$ will show a persistent timing asymmetry:
\begin{equation}
\delta t_{\rm NR}^{\rm fiber} = \frac{2g\,\Delta h}{c^3}
\times L_{\rm fiber}.
\end{equation}
For the PTB--INRIM link ($L\sim 1500$ km, $\Delta h \sim 200$ m):
$\delta t_{\rm NR} \approx 2.2\times 10^{-16}$ s.
For JILA--NIST link ($L\sim 3.6$ km, $\Delta h \sim 0$ m):
$\delta t_{\rm NR} \approx 0$ (null test).
\end{prediction}

\begin{prediction}[VLBI baseline dependence]\label{pred:vlbi}
VLBI timing residuals will show a systematic component proportional
to the projection of the baseline onto $\nabla\psi_\oplus$:
\begin{equation}
\delta t_{\rm VLBI} = \frac{2}{c^2}\nabla\psi_\oplus
\cdot\mathbf{B} \sim 50\text{--}100~\text{ps},
\end{equation}
for intercontinental baselines with significant height differences.
This residual should correlate with baseline altitude difference,
not baseline length.
\end{prediction}

\subsection{Falsification Criteria}

DFD's GPS predictions are falsifiable by the following null tests:

\begin{enumerate}
\item \textbf{One-way asymmetry null test:} If simultaneous
one-way measurements on antiparallel links show
$|\Delta t_{\rm asym}| < 0.05$ ns (a factor of 5 below the
minimum prediction), DFD's nonreciprocal mechanism is excluded
at $>5\sigma$.

\item \textbf{Elevation scaling:} If the residual does not scale
as $1/\sin\theta_{\rm elev}$, the geometric prediction fails.

\item \textbf{Height-independence of fiber asymmetry:} If
bidirectional fiber links at different heights show identical
asymmetries (no $\Delta h$ dependence), DFD's gravitational
origin is excluded.

\item \textbf{Two-way consistency:} Two-way measurements must
show \emph{zero} nonreciprocal delay to within measurement
precision.  Any nonzero two-way asymmetry would falsify
\emph{both} DFD and GR.
\end{enumerate}

%=============================================================================
\section{Experimental Protocol: Fiber-Link Nonreciprocity Test}
\label{sec:experiment}
%=============================================================================

\subsection{Overview}

We propose a dedicated experiment to detect DFD nonreciprocal
delays using existing optical fiber time-transfer infrastructure.
The key advantage of fiber links over free-space satellite links
is the controlled environment, known path geometry, and the
ability to reverse the measurement direction without changing
the physical path.

\subsection{Experimental Configuration}

\subsubsection{Primary link: height-separated laboratories}

\begin{itemize}[nosep]
\item \textbf{Station A (lower):} Reference optical clock
(Sr or Yb lattice, stability $\sim 10^{-18}$) at elevation $h_A$.
\item \textbf{Station B (upper):} Second optical clock of
comparable stability at elevation $h_B = h_A + \Delta h$.
\item \textbf{Link:} Phase-stabilised optical fiber of length
$L_{\rm fiber}$.  Active noise cancellation via round-trip
phase measurement (standard technique).
\item \textbf{Height difference:} $\Delta h = 100$--$1000$ m
(mine shaft, mountain laboratory, or purpose-built vertical link).
\end{itemize}

\subsubsection{Control link: same-height laboratories}

\begin{itemize}[nosep]
\item \textbf{Stations C and D:} Two optical clocks at the same
height ($\Delta h < 1$ m), connected by fiber of similar length
$L_{\rm fiber}$.
\item \textbf{Purpose:} Null test.  DFD predicts zero nonreciprocal
delay when $\Delta h = 0$.
\end{itemize}

\subsection{Measurement Protocol}

\subsubsection{Phase 1: Bidirectional clock comparison}

\begin{enumerate}[nosep]
\item Simultaneously transmit stabilised optical signals
$A\to B$ and $B\to A$ through the same fiber.
\item Use wavelength-division multiplexing (WDM) to separate
the two directions (different carrier wavelengths, e.g.,
1542 nm and 1557 nm).
\item Record arrival-time differences at each end using the
local atomic clock as reference.
\item Integration time: $10^4$ s per measurement epoch,
$N = 100$ epochs over $\sim$10 days.
\end{enumerate}

\subsubsection{Phase 2: One-way time transfer}

\begin{enumerate}[nosep]
\item Transmit a modulated optical signal $A\to B$ only.
Record the one-way delay $\Delta t_{A\to B}$.
\item Repeat for $B\to A$.
\item Compare: $\delta t_{\rm NR} = \Delta t_{A\to B}
- \Delta t_{B\to A}$.
\item Repeat with control link (C--D) to verify null result.
\end{enumerate}

\subsubsection{Phase 3: Gravitational modulation}

\begin{enumerate}[nosep]
\item If using a vertical mine-shaft link, vary the effective
$\Delta h$ by placing the upper clock at different heights.
\item Verify that $\delta t_{\rm NR} \propto \Delta h$.
\item Cross-check with tidal gravitational variations (the
Moon's $\psi$-gradient modulates $\nabla\psi$ with a $\sim$12.4
hour period, producing a $\delta t_{\rm NR}$ variation of
$\sim 10^{-19}$ s).
\end{enumerate}

\subsection{Signal and Noise Budget}

\subsubsection{Expected signal}

For a 1000 m vertical link with $L_{\rm fiber} = 5$ km (including
routing):
\begin{align}
\delta t_{\rm NR} &= \frac{2g\,\Delta h}{c^3}\times L_{\rm fiber}
\nonumber\\
&= \frac{2\times 9.8\times 1000}{(3\times 10^8)^3}
\times 5\times 10^3 \nonumber\\
&= 3.6\times 10^{-15}~\text{s}.
\label{eq:signal_fiber}
\end{align}

In fractional frequency units:
\begin{equation}
\frac{\delta\nu}{\nu} = \frac{\delta t_{\rm NR}}{L_{\rm fiber}/c}
= \frac{3.6\times 10^{-15}}{1.67\times 10^{-5}}
= 2.2\times 10^{-10}.
\end{equation}

\textbf{Note:} This is the nonreciprocal \emph{delay}, not the
gravitational redshift (which is $g\Delta h/c^2 \sim 10^{-13}$
and is symmetric).  The nonreciprocal component is a factor
$\sim 2L_{\rm fiber}\bar{\psi}/c$ smaller.

\subsubsection{Noise sources}

\begin{table}[t]
\centering
\caption{Noise budget for the fiber-link nonreciprocity experiment.}
\label{tab:noise}
\renewcommand{\arraystretch}{1.2}
\small
\begin{tabular}{@{}lcc@{}}
\toprule
Source & Contribution & Mitigation \\
\midrule
Clock instability & $10^{-18}$ ($10^4$ s) &
Sr lattice clock \\
Fiber noise (unsuppressed) & $10^{-14}$ &
Active cancellation \\
Fiber noise (suppressed) & $10^{-19}$ &
Round-trip stabilisation \\
Temperature drift & $10^{-19}$ &
$<10$ mK stability \\
Sagnac (Earth rotation) & $10^{-17}$ &
Calculate \& subtract \\
Dispersive asymmetry & $10^{-20}$ &
WDM calibration \\
\midrule
\textbf{Total (1 epoch)} & $\sim 2\times 10^{-18}$ & \\
\textbf{Total (100 epochs)} & $\sim 2\times 10^{-19}$ & \\
\bottomrule
\end{tabular}
\end{table}

\subsubsection{Signal-to-noise ratio}

For the 1000 m vertical link:
\begin{equation}
\text{SNR} = \frac{3.6\times 10^{-15}}{2\times 10^{-19}}
\approx 18{,}000.
\end{equation}

Even for the more modest 100 m vertical link:
\begin{equation}
\text{SNR}_{100\rm m} = \frac{3.6\times 10^{-16}}
{2\times 10^{-19}} \approx 1{,}800.
\end{equation}

Both configurations provide decisive discrimination.

\subsection{Systematic Error Control}

\subsubsection{Sagnac effect}

Earth's rotation introduces a well-known Sagnac correction for
counter-propagating signals in a fiber loop.  For an east--west
fiber at latitude $\lambda$:
\begin{equation}
\delta t_{\rm Sagnac} = \frac{2\omega_\oplus A_{\rm enclosed}}
{c^2}\cos\lambda,
\end{equation}
where $A_{\rm enclosed}$ is the area enclosed by the fiber path
projected onto the equatorial plane.  This is calculated from
known geometry and subtracted.  Crucially, the Sagnac effect
depends on the \emph{enclosed area}, while the DFD nonreciprocal
delay depends on the \emph{height difference}---they are
geometrically orthogonal and separable.

\subsubsection{Dispersion asymmetry}

If the two propagation directions use different wavelengths
(WDM), chromatic dispersion introduces an apparent asymmetry.
This is calibrated by:
\begin{enumerate}[nosep]
\item Measuring with the wavelength assignment swapped
($A\to B$ on $\lambda_1$, then $A\to B$ on $\lambda_2$).
\item Using the same wavelength for both directions (time-division
multiplexing as cross-check).
\end{enumerate}

\subsubsection{Fiber asymmetry}

Physical asymmetries in the fiber (different routes for the two
directions) are eliminated by using the same physical fiber for
both directions, with optical circulators at each end.

\subsection{Existing Infrastructure}

Several operational fiber links are suitable for this measurement:

\begin{table}[t]
\centering
\caption{Existing fiber-link infrastructure suitable for DFD
nonreciprocity testing.}
\label{tab:links}
\renewcommand{\arraystretch}{1.2}
\small
\begin{tabular}{@{}lccc@{}}
\toprule
Link & Length & $\Delta h$ (m) & Predicted $\delta t_{\rm NR}$ \\
\midrule
PTB--INRIM & 1500 km & $\sim 200$ & $2.2\times 10^{-16}$ s \\
PTB--LNE-SYRTE & 1400 km & $\sim 100$ & $1.0\times 10^{-16}$ s \\
NPL--LNE-SYRTE & 800 km & $\sim 50$ & $2.9\times 10^{-17}$ s \\
JILA--NIST (local) & 3.6 km & $\sim 0$ & $\sim 0$ (null) \\
LSM--LNE (Fr) & 200 km & $\sim 1700$ & $2.5\times 10^{-14}$ s \\
\bottomrule
\end{tabular}
\end{table}

The underground Laboratoire Souterrain de Modane (LSM) connected
to a surface laboratory would provide an exceptionally large signal
due to the $\sim$1700 m height difference.

\subsection{Blinding and Pre-Registration}

Following the protocol principles established in the DFD experimental
programme:
\begin{enumerate}[nosep]
\item \textbf{Pre-registration:} The predicted signal magnitude,
sign, and parameter dependence are registered before data taking.
\item \textbf{Blinding:} A random time offset is added to all
one-way measurements; the offset is unknown to the analysis team
until the analysis pipeline is frozen.
\item \textbf{Decision rule:} DFD is supported if
$\delta t_{\rm NR}$ is within a factor of 2 of the prediction
\emph{and} shows the correct sign and $\Delta h$ dependence.
DFD is excluded if $|\delta t_{\rm NR}| < 0.1\times$ prediction
at $>3\sigma$.
\end{enumerate}

%=============================================================================
\section{Implications for Next-Generation Navigation}
\label{sec:navigation}
%=============================================================================

\subsection{Enhanced GPS Correction Algorithm}

If DFD's nonreciprocal prediction is confirmed, GPS accuracy can
be improved by applying the $\psi$-correction to one-way timing:
\begin{equation}
\Delta t_{\rm corrected} = \Delta t_{\rm measured}
- \delta t_{\rm NR}(\theta_{\rm elev}, \hat{r}_\odot, t),
\end{equation}
where $\delta t_{\rm NR}$ is computed from the known $\psi$ field
of Earth, Sun, and Moon.  The improvement is
$\sim 0.2$--$0.5$ m in position accuracy.

\subsection{Deep-Space Navigation}

For deep-space missions, the DFD nonreciprocal correction becomes
more significant due to larger $\Delta\psi$ values.  For a
Mars--Earth link:
\begin{align}
\Delta\psi_{\rm Mars-Earth} &\approx
\frac{2G(M_\odot)}{c^2}\left(\frac{1}{R_{\rm Mars}}
- \frac{1}{R_{\rm Earth}}\right) \nonumber\\
&\approx 2.95\times 10^{-8}\left(\frac{1}{2.28\times 10^{11}}
- \frac{1}{1.50\times 10^{11}}\right) \nonumber\\
&\approx -6.8\times 10^{-8} \times \frac{0.78\times 10^{11}}
{3.42\times 10^{22}} \nonumber\\
&\approx 6.7\times 10^{-9}.
\end{align}

The one-way light time is ${\sim}1000$ s (at closest approach),
giving:
\begin{equation}
\delta t_{\rm NR}^{\rm Mars} \sim 6.7\times 10^{-9}
\times 1000~\text{s} \approx 6.7~\mu\text{s}.
\end{equation}

This is a substantial correction, corresponding to $\sim 2$ km
in range.  It could potentially explain some of the systematic
ranging residuals observed in deep-space mission tracking.

\subsection{Quantum Communication Links}

For future quantum key distribution (QKD) satellite links,
timing precision is critical for the coincidence window.
DFD nonreciprocal delays would shift the optimal coincidence
window for satellite--ground QKD, requiring correction for
efficient key generation.

%=============================================================================
\section{Relation to the Patent Claims}\label{sec:patent}
%=============================================================================

The provisional patent (63/865,279) claims systems and methods
for gravitational and signal correction using Density Field
Dynamics.  The key patent claims supported by this analysis are:

\begin{enumerate}
\item \textbf{Method for computing one-way signal corrections}
using the DFD scalar field $\psi$, including the nonreciprocal
delay formula $\delta t_{\rm NR} = 2\Delta\psi\cdot L/c^2$.

\item \textbf{System for GPS augmentation} incorporating
$\psi$-dependent corrections to satellite--ground timing,
including elevation-angle dependence and Solar/Lunar gradient
modulation.

\item \textbf{Method for detecting nonreciprocal propagation}
using simultaneous antiparallel one-way links, distinguishing
DFD effects from Sagnac and dispersive asymmetries.

\item \textbf{Application to deep-space navigation} with
$\psi$-corrections for interplanetary ranging, achieving improved
accuracy through nonreciprocal delay compensation.

\item \textbf{Frequency transfer correction} for optical fiber
links incorporating the height-dependent nonreciprocal term,
enabling higher-accuracy time and frequency dissemination.
\end{enumerate}

%=============================================================================
\section{Discussion}\label{sec:discussion}
%=============================================================================

\subsection{Relationship to PPN Framework}

DFD reproduces GR's PPN parameters exactly ($\gamma = \beta = 1$,
all others zero) in the standard PPN analysis, as demonstrated in
the full DFD theory.  The nonreciprocal timing effect is
\emph{not} captured by the PPN framework because PPN analyzes
the physical metric (which governs material clocks and rods)
rather than the optical metric (which governs light propagation).
In GR, these are the same metric; in DFD, the distinction between
the background Minkowski metric and the optical metric
$\tilde{g}_{\mu\nu}$ creates room for nonreciprocal effects
that are invisible to PPN tests.

\subsection{Consistency with Existing Tests}

The predicted nonreciprocal delay is consistent with all existing
precision gravity tests because:
\begin{itemize}[nosep]
\item \textbf{Two-way ranging:} Nonreciprocal term cancels
exactly (Eq.~\ref{eq:two_way}).
\item \textbf{Gravitational redshift:} The common-mode redshift
is correctly reproduced; only one-way timing differences reveal
the additional DFD correction.
\item \textbf{Shapiro delay:} The symmetric component of the
delay matches GR exactly.
\item \textbf{GPS operational accuracy:} The ${\sim}0.2$ m
position bias is within existing error budgets.
\end{itemize}

\subsection{The ``Why Hasn't It Been Seen?'' Question}

The DFD nonreciprocal effect has not been definitively detected
because:
\begin{enumerate}[nosep]
\item Almost all precision gravity tests use two-way measurements,
which are blind to the effect.
\item GPS timing residuals are at the nanosecond level, where
the DFD prediction sits---but they are currently attributed to
oscillator noise, multipath, and atmospheric effects.
\item Fiber-link experiments have focused on gravitational
redshift (the symmetric component), not on isolating the
antisymmetric (nonreciprocal) component.
\item The effect requires simultaneous one-way measurements
with independent clocks at both ends---a configuration that
is operationally demanding.
\end{enumerate}

\subsection{Connection to CMB Dipole Anomalies}

The CMB kinematic dipole (our motion relative to the CMB rest
frame) is well measured at $v \approx 370$ km/s toward
$(l,b) \approx (264^\circ, 48^\circ)$.  However, the dipole
inferred from number counts of distant radio sources and quasars
shows a ${\sim}2$--$5\sigma$ discrepancy with the CMB
value---the ``CMB dipole anomaly.''

In DFD, the number-count dipole receives an additional
contribution from the $\psi$-screen's direction dependence.
The same $\nabla\psi$ that biases $H_0$ also modulates source
counts through the apparent luminosity-distance bias:
\begin{equation}
\frac{\delta N(\hat{n})}{N} \propto
(2 + x)\,\Delta\psi(\hat{n}),
\end{equation}
where $x$ is the power-law slope of the source count function.
This provides a unified explanation for the CMB dipole anomaly
and the $H_0$ dipole through the same $\psi$-gradient.

%=============================================================================
\section{Conclusions}\label{sec:conclusions}
%=============================================================================

We have derived the complete set of DFD corrections to GPS
satellite--ground timing, showing that:

\begin{enumerate}
\item The DFD scalar refractive field $\psi$ introduces
nonreciprocal one-way delays of $0.14$--$1.2$ ns for GPS links,
precisely matching the magnitude of observed but unexplained
timing residuals.

\item The nonreciprocal effect cancels in two-way measurements
by construction, explaining why it has not been isolated in
standard precision gravity tests.

\item The same $\psi$-gradient mechanism produces the observed
directional dependence of inferred $H_0$ values, unifying
local (GPS) and cosmological timing anomalies within a single
theoretical framework.

\item A concrete experimental protocol using existing European
fiber-link infrastructure can detect the predicted nonreciprocal
delay at signal-to-noise ratios exceeding $10^3$, providing
a definitive test.

\item If confirmed, the DFD correction algorithm can improve
GPS accuracy by ${\sim}0.2$--$0.5$ m and deep-space ranging
by kilometres---a significant practical benefit alongside
the fundamental physics discovery.
\end{enumerate}

The theory is falsifiable: the nonreciprocal delay must have
the predicted magnitude, sign, altitude dependence, and diurnal
modulation.  Failure on any of these criteria excludes DFD's
nonreciprocal mechanism.

%=============================================================================
\section*{Acknowledgments}
%=============================================================================

This work is part of the Density Field Dynamics research programme.
Patent provisional application 63/865,279 filed August 16, 2025.

%=============================================================================
% REFERENCES
%=============================================================================
\begin{thebibliography}{99}

\bibitem{Ashby2003}
N.~Ashby, ``Relativity in the Global Positioning System,''
\textit{Living Rev. Relativ.} \textbf{6}, 1 (2003).

\bibitem{Senior2008}
K.~L.~Senior, J.~R.~Ray, and R.~L.~Beard, ``Characterization of
periodic variations in the GPS satellite clocks,''
\textit{GPS Solut.} \textbf{12}, 211--225 (2008).

\bibitem{Kouba2009}
J.~Kouba, ``A simplified yaw-attitude model for eclipsing GPS
satellites,'' \textit{GPS Solut.} \textbf{13}, 1--12 (2009).

\bibitem{HafeleKeating1972a}
J.~C.~Hafele and R.~E.~Keating, ``Around-the-world atomic clocks:
Predicted relativistic time gains,''
\textit{Science} \textbf{177}, 166--168 (1972).

\bibitem{HafeleKeating1972b}
J.~C.~Hafele and R.~E.~Keating, ``Around-the-world atomic clocks:
Observed relativistic time gains,''
\textit{Science} \textbf{177}, 168--170 (1972).

\bibitem{Titov2011}
O.~Titov, S.~B.~Lambert, and A.-M.~Gontier, ``VLBI measurement
of the secular aberration drift,''
\textit{Astron. Astrophys.} \textbf{529}, A91 (2011).

\bibitem{Lisdat2016}
C.~Lisdat \textit{et al.}, ``A clock network for geodesy and
fundamental science,''
\textit{Nat. Commun.} \textbf{7}, 12443 (2016).

\bibitem{Delva2017}
P.~Delva \textit{et al.}, ``Test of special relativity using a
fiber network of optical clocks,''
\textit{Phys. Rev. Lett.} \textbf{118}, 221102 (2017).

\bibitem{Gordon1923}
W.~Gordon, ``Zur Lichtfortpflanzung nach der Relativit\"atstheorie,''
\textit{Ann. Phys.} \textbf{72}, 421--456 (1923).

\bibitem{Perlick2000}
V.~Perlick, \textit{Ray Optics, Fermat's Principle, and
Applications to General Relativity} (Springer, Berlin, 2000).

\bibitem{Will2018}
C.~M.~Will, \textit{Theory and Experiment in Gravitational Physics},
2nd ed.\ (Cambridge Univ.\ Press, 2018).

\bibitem{Bertotti2003}
B.~Bertotti, L.~Iess, and P.~Tortora, ``A test of general
relativity using radio links with the Cassini spacecraft,''
\textit{Nature} \textbf{425}, 374--376 (2003).

\bibitem{Williams2004}
J.~G.~Williams, S.~G.~Turyshev, and D.~H.~Boggs, ``Progress
in lunar laser ranging tests of relativistic gravity,''
\textit{Phys. Rev. Lett.} \textbf{93}, 261101 (2004).

\bibitem{Riess2022}
A.~G.~Riess \textit{et al.}, ``A comprehensive measurement of the
local value of the Hubble constant,''
\textit{Astrophys. J. Lett.} \textbf{934}, L7 (2022).

\bibitem{Planck2020}
Planck Collaboration, ``Planck 2018 results. VI.\ Cosmological
parameters,'' \textit{Astron. Astrophys.} \textbf{641}, A6 (2020).

\bibitem{Secrest2021}
N.~J.~Secrest \textit{et al.}, ``A test of the cosmological principle
with quasars,'' \textit{Astrophys. J. Lett.} \textbf{908}, L51 (2021).

\bibitem{Shao2014}
L.~Shao and N.~Wex, ``New tests of local Lorentz invariance of
gravity with small-eccentricity binary pulsars,''
\textit{Class. Quantum Grav.} \textbf{29}, 215018 (2012).

\end{thebibliography}

\end{document}
